\documentclass[10pt,a4paper]{amsart}
\usepackage[margin=19mm]{geometry}
\usepackage{amsmath,amssymb,mathtools}
\usepackage[scr=rsfs,cal=euler]{mathalfa}
\usepackage{xfrac}
\usepackage[unicode,hidelinks,bookmarks=false]{hyperref}
\newcommand{\ie}{i.e.\ }
\newcommand{\cf}{cf.\ }
\newcommand{\resp}{respectively\ }
\newcommand{\Subsection}{\subsection}
\providecommand{\nobreakdash}{\nobreak\hbox{-}}
\setlength{\emergencystretch}{2em}
\multlinegap0pt
\usepackage{amssymb}
\newtheorem{lemma}{Lemma}
\title{Asymptotics of the coefficients of polynomials, asymptotic zero distribution and free probability}
\author{Walter Van Assche}
\date{Source: arXiv:2609.12717v1 (CC BY 4.0)}
\begin{document}
\maketitle

\noindent\emph{Source and licence.} Asymptotics of the coefficients of polynomials, asymptotic zero distribution and free probability. Version arXiv:2609.12717v1 (CC BY 4.0), distributed by arXiv under the Creative Commons Attribution 4.0 International licence, http://creativecommons.org/licenses/by/4.0/, as recorded in the arXiv licence metadata for this exact version and shown on the abstract page https://arxiv.org/abs/2609.12717v1. Full licence notice: \texttt{licenses/arxiv-2609.12717-CC-BY-4.0.txt}.\par\smallskip

\noindent\textit{Editorial scope.} Counted unit: the article's lemma lemma (source0.tex lines 402--406). The article's complete original proof of it is delivered with it (source0.tex lines 408--438, one \texttt{proof} environment of 31 lines). No mathematics is written, repaired or re-derived: the statement and the proof are the article's own lines, in the article's own notation, and no retained line is substituted or re-flowed.  Retained uncounted as the source-local context the proof needs: lines 397--399.  Nothing was omitted from any retained span, so the package carries no omission record.  The retained lines cite 1 works, whose entries are transcribed verbatim from the article's own bibliography source in the e-print and delivered with short editorial labels recorded in the item's reference mapping.  No retained line contains a figure environment, an image-inclusion command or drawing code, so no image is delivered and the package declares no asset dependency.  Editorial formatting only, in addition: an AMS \texttt{amsart} preamble that restates the article's own declarations and macros used by the retained lines, namely the environments \texttt{lemma} on separate counters, together with the standard packages those lines need.  The retained environments print in this document as Lemma 1 in order of appearance, whereas the article prints the same items with the numbering of its own full document.  The complete native source of the article as distributed in the arXiv e-print for this exact version is bundled unchanged as \texttt{sources/210142/source0.tex}.
\begin{equation}  \label{coefCharlier}
   {}_2F_0\left( \begin{array}{c}   -k, \nu+n-k+1 \\  - \end{array}; -c  \right) = C_k(-\nu-n+k-1;1/c) .
\end{equation}

\begin{lemma}
The following limit holds
\[    \lim_{k/n \to t}   \frac{1}{n}  \frac{{}_2F_0\left( \begin{array}{c}   -k, \nu+n-k+1 \\  - \end{array}; -c  \right)}
                                                               { {}_2F_0\left( \begin{array}{c}   -k+1, \nu+n-k+2 \\  - \end{array}; -c  \right)}  = c (1-t).    \]
\end{lemma}

\begin{proof}
Let
\[   f(n,k) =  {}_2F_0\left( \begin{array}{c}   -k, \nu+n-k+1 \\  - \end{array}; -c  \right), \]
then one has the recurrence
\begin{equation}   \label{f(n,k)rec}
    f(n,k) + (2ck-cn-c\nu-2c-1) f(n,k-1) + c^2(k-1)(-\nu-n+k-2) f(n,k-2).  
\end{equation}
This can be obtained by using the recurrence relation for Charlier polynomials, combined with the difference relations for Charlier polynomials,
or simply by using the \texttt{hypersum} command in the maple package \texttt{sumtools}. Here is a proof using Charlier polynomials.
The three term recurrence relation for the Charlier polynomials is \cite[Eq (9.14.3) on p.~247]{koekoek}
\begin{equation}  \label{Charlier3TR}
  -x C_n(x;a) = a C_{n+1}(x;a) - (n+a) C_n(x;a) + n C_{n-1}(x;a).   
\end{equation}
Take $n=k$, $a=1/c$ and $x=-\nu-n+k-1$ to find
\begin{multline*}  c(\nu+n-k+1) C_k(-\nu-n+k-1;1/c) \\
= C_{k+1}(-\nu-n+k-1;1/c) - (kc+1) C_k(-\nu-n+k-1;1/c) + kc C_{k-1}(-\nu-n+k-1;1/c). 
\end{multline*}
For  $C_{k+1}$ we use the backward shift operator \cite[Eq. (9.14.8) on p.~248]{koekoek} 
\[      C_{k+1}(-\nu-n+k;1/c) - C_{k+1}(-\nu-n+k-1;1/c) = -(k+1) c C_k(-\nu-n+k-1;1/c),  \]
and for $C_{k-1}$ we use the forward shift operator \cite[Eq. (9.14.6) on p.~248]{koekoek} 
\[     C_{k-1}(-\nu-n+k-1;1/c) - c(-\nu-n+k-1) C_{k-1}(-\nu-n+k-2;1/c) = C_k(-\nu-n+k-1;1/c).  \]
Together with the relation \eqref{coefCharlier} we then find the recurrence \eqref{f(n,k)rec}.

Divide the recurrence \eqref{f(n,k)rec} by $n f(n,k-1)$, then the limit $L(t) = \lim_{k/n \to t} \frac{1}{n} \frac{f(n,k)}{f(n,k-1)}$ satisfies
\[  L(t) + c(2t-1) - c^2t(1-t) \frac{1}{L(t)} = 0.  \]
Solving this gives
\[   L(t) = \frac{-c(2t-1) \pm c}{2}.  \]
Note that $f(n,k)$ is a positive quantity since
\[     f(n,k) = \sum_{j=0}^k \binom{k}{j} (\nu+n-k+1)_j c^j, \]
so that we need to choose the positive solution, giving $L(t) = c(1-t)$.
\end{proof}

\begin{thebibliography}{99}
\bibitem[koekoe]{koekoek} Roelof Koekoek, Peter A. Lesky, Ren\'e F. Swarttouw, \textit{Hypergeometric Orthogonal Polynomials and Their $q$-Analogues}, Springer Monographs in Mathematics, Springer-Verlag, Berlin Heidelberg, 2010.
\end{thebibliography}
\end{document}
